\paragraph{Related work}

Ling et al.~\cite{ling_table2latex_2025} propose a method for \LaTeX{} code generation from table images using reinforced multimodal language models. Their work highlights the potential of leveraging advanced AI techniques to improve the accuracy and efficiency of \LaTeX{} document preparation. However, as they claim, the method is suffering from low-efficiency and high token cost due to the compilation principles of \LaTeX{}.

Paster et al.~\cite{paster2023openwebmath} attach the importance to mathematical content for LLM training. They constructed a 14.7B token dataset from web sources and demonstrated the complexity of extracting and processing \LaTeX{} content from diverse formats.

Kale and Nadadur~\cite{kale2025texpert} demonstrate that even state-of-the-art LLMs struggle with \LaTeX{} code generation. The LLMs achieved only 15\% accuracy on complex tasks, with logical errors accounting for 54\% of failures.

Xia et al.~\cite{xia2024docgenome} performed large-scale benchmarks to reveal that Equation-to-\LaTeX{} and Table-to-\LaTeX{} conversion tasks remain challenging even with extensive training data (Edit Distance > 0.21). The reason is the high information entropy of \LaTeX{} source code.

Lin et al.~\cite{lin2024accurate} highlight that the high computational cost of applying LLMs directly to document processing. It achieves up to 30$\times$ cost reduction through semantic hierarchical indexing. They launch an appeal for the need of more efficient document representations.

{\color{red}
\textbf{The \TeX{}/\LaTeX{} ecosystem and its evolution.}
The \TeX{} system, originally designed by Knuth~\cite{knuth1984texbook} in the late 1970s, was later extended by Lamport~\cite{lamport1994document} into \LaTeX{}, which became the de facto standard for scientific typesetting. Mittelbach et al.~\cite{mittelbach2004latex} provided a comprehensive treatment of \LaTeX{}'s package ecosystem and its growing complexity. The \LaTeX{}3 project, initiated in the 1990s, attempted to address engineering deficiencies through the \texttt{expl3} programming layer, which was formally integrated into the \LaTeX{} kernel in 2020~\cite{latex3team2024expl3}. However, this progressive integration has not resolved the fundamental structural contradictions of \LaTeX{}2$\varepsilon$. Berry et al.~\cite{berry2025texlive} document the continuous expansion of TeX~Live, whose ISO size has grown from 2.4\,GB in 2008 to 5.9\,GB in 2025, reflecting the ever-accumulating complexity of the distribution model.

\textbf{Alternative typesetting and structured editing systems.}
The tension between batch-compiled markup languages and interactive editing has motivated several alternative systems. GNU \TeX{}macs~\cite{van_der_hoeven_gnu_2001} pioneered the WYSIWYG structured editing paradigm, demonstrating that high-quality mathematical typesetting is achievable without \TeX{}'s compilation model. Typst~\cite{madje2023typst}, introduced in 2023, offers a modern markup-based typesetting system with incremental compilation, native Unicode support, and a scripting layer, attracting significant community adoption. LyX~\cite{lyx2024} provides a quasi-WYSIWYG front-end to \LaTeX{} but remains fundamentally dependent on the \TeX{} compilation pipeline. Overleaf~\cite{overleaf_docs} has emerged as the dominant cloud-based \LaTeX{} collaboration platform, but its architecture still relies entirely on traditional \LaTeX{} compilation, inheriting all associated limitations in preview latency, conflict resolution, and extensibility.

\textbf{LLMs for scientific document processing.}
The integration of LLMs into scientific writing workflows has exposed the token inefficiency and syntactic complexity of \LaTeX{}. Xia et al.~\cite{xia2024overleafcopilot} explored the integration of LLMs with Overleaf, demonstrating both the potential and the constraints imposed by the underlying \LaTeX{} compilation model. Filho~\cite{filho_doa_2026} proposed documentation-oriented architectures using Markdown as a coordination layer in multi-agent ecosystems, highlighting the limitations of verbose markup languages for AI-driven workflows. Wang et al.~\cite{wang2024paperqa} developed PaperQA for scientific document question answering, where the complexity of parsing \LaTeX{} source remains a bottleneck. The broader trend toward AI-assisted scientific discovery~\cite{lu2024aiscientist} further underscores the need for document formats with lower information entropy and higher structural transparency.
}
